\section{R4. Indexed families and the motive problem}\label{sec:dependent}

\subsection{The problem and corrected starting point}\label{sec:dependent-scope}

The historical evidence distinguishes elimination from synthesis of a recursive
motive. Head of a nonempty vector succeeds through native dependent
\code{cases}; vector map requires
\begin{lstlisting}
(A B : Type) -> (n : Nat) -> (A -> B) -> Vec A n -> Vec B n
\end{lstlisting}
with motive $M(n,x)=\mathsf{Vec}\ B\ n$. The inspected structural-recursion
rule excluded indexed families and \code{Nat} before invoking Lean induction.
Those source restrictions are not limitations of Lean's induction facilities.
An initial supported indexed adapter is consequently a separate, smaller
milestone than unrestricted motive invention.\prov{N1--N9}

Four capabilities remain distinct: infer a motive for a prepared target;
generalize locals so recursive hypotheses apply at changing arguments; repair
propositional index equalities with transport; and expose a type family
$F(?n)$ whose head depends on an unknown index. General higher-order unification
is undecidable, while pattern unification has decidable fragments and
practical elaborators add postponement and pruning\cite{huet75,miller91,abelpientka11}.
Native dependent pattern matching supplies index unification, not a complete
enumerator of all useful strengthened motives\cite{cockx}.

Canonical already demonstrates indexed recursor applications and generalized
induction with a supplied stronger principle and predicate synthesis. The
former claim that existing search never handles indexed or generalized
induction is therefore too strong. Its translation restrictions remain
relevant, and neither example establishes arbitrary motive invention.
Agsy/Mimer, Idris proof search, and \code{sauto} remain useful comparisons;
\code{sauto}'s dependent elimination and unfolding should not be mistaken for
automatic induction discovery, which its documentation explicitly excludes.
\cite{canonical,agsy,sauto,sautodocs}\prov{N1,N3--N9}

\subsection{A finite motive language with explicit generalization}\label{sec:dependent-motives}

For a major premise $x:I\,\vec p\,\vec\imath$, separate three choices:
\begin{enumerate}
\item abstract the indices and major premise already present in the target;
\item revert a dependency-closed telescope of existing parameters;
\item invent a bounded carrier, auxiliary result, or invariant absent from the
original context.
\end{enumerate}
Mandatory dependency generalization is primarily an elaboration task;
optional reversion is a search choice; invention is a larger synthesis problem.
They share the carrier service of Section~\ref{sec:carriers}, but should not
be charged to one undifferentiated powerset search.\prov{N1--N9}

Fix scoped goal syntax and a finite collection of eligible occurrence positions.
With $k$ positions, at most $2^k$ raw subsets arise before scope, typing,
symmetry, and justified equality filters. This proves finite syntactic
generation, not low complexity or completeness for higher-order unification.
Nonoverlapping maximal occurrences are one useful policy, not the only one:
a target containing $\mathsf{Vec}\ B\ (n+k)$ may need abstraction of the
occurrence of $n$ inside the index while retaining $k$. A policy abstracting
only the entire maximal $n+k$ need not generate that motive. Occurrences are
scoped expressions with dependencies, not printed substrings.\prov{N3,N4}

Nor can occurrence abstraction invent a new binder or result component:
\[
 M(n,x)=\mathsf{Vec}\ B\ n,\qquad
 M(n,x)=\prod_{a:S(n,x)}T(n,x,a),\qquad
 M(n,x)=T(n,x)\times U(n,x)
\]
are different recipe families. The latter two can introduce $S$, $U$, or a
strengthened invariant absent from the original target. Explicitly close a
bounded template language under dependent products, products, dependent pairs,
subtypes, option-like initialization state, and permitted carrier constructors.
Fix its terminals, literals, universe-expression size, auxiliary type formers,
number of new binders, and total expression size; otherwise a nominal syntax
bound need not produce a finite space.\prov{N1--N9}

Generalization preserves telescope order and closes over both local declaration
types and values. If $j:\mathsf{Fin}\ n$, moving $n$ across induction while
leaving $j$ fixed is ill scoped. If a recursive call changes accumulator $a$,
generalize $a$ even when its type mentions no index. Conversely, reverting
every local makes branches unnecessarily large. For vector map, keep $A$, $B$,
and $f:A\to B$ fixed; the empty branch produces $\mathsf{Vec}\ B\ 0$ and the
cons branch combines $f$ of the head with the induction hypothesis at the
tail length. Vector zip and lookup require the stronger parameter interfaces
already traced in the unified proposal, not independent rediscovery here.

Compound or repeated indices need explicit preparation. Replace a compound
index such as $n+m$ by a fresh variable only while retaining its defining
equality and dependent consequences. Then use native elimination to create
the obligations. A generalized variable without that equation changes the
problem. A motive producing proof-carrying results must retain dependence on
the major value itself as well as its numerical indices.\prov{N1--N9}

\Needspace{20\baselineskip}
\paragraph{The recipe algorithm.}\label{sec:dependent-recipe}
The bounded service enumerates scoped recipes and validates each native
eliminator attempt before exposing its branches.
\begin{lstlisting}[caption={Bounded motive service; pseudocode.}]
enumerateMotiveRecipes(goal, major, bounds):
  expose the inductive head and read native recursor metadata
  prepare mandatory dependent telescope and index equations
  yield native inferred motive
  yield small dependency-closed generalizations
        ranked by failed recursive-hypothesis applications
  yield bounded parameter/carrier/result strengthenings
  yield the declared wider scoped-occurrence abstractions
  for each recipe, within its fair bounded schedule:
    validate scope, sort, universes, and binder dependencies
    attempt native eliminator construction transactionally
    retain branches, substitutions, transports, and replay recipe
    check final specialization back to the original target
\end{lstlisting}
A recursive hypothesis fixed at the wrong accumulator suggests generalizing
that argument and its dependents. A missing result component suggests product
strengthening; an algebraic reassociation may instead need a lemma.
These explanations rank proposals, not prove that omitted recipes cannot work.
Prefer the native inferred motive before expensive alternatives and retain a
fair bounded fallback.\prov{N1--N9}

\begin{proposition}[Restricted motive completeness]\label{prop:dependent-motives}
Fix a finite scoped occurrence set, a finite family of dependency-closed
generalization telescopes, and a finite result-lifting grammar. If candidate
construction and exact checking terminate, enumeration of every combination,
discarding only ill-typed combinations, terminates and finds every well-typed
motive in that class.
\end{proposition}
\begin{proof}
There are finitely many combinations and each is processed in finite time.
Every member of the specified class is generated by definition, and every
well-typed member survives the stated filter.
\end{proof}
A timeout is not exact checking and remains unresolved. This theorem describes
exhaustion of the declared recipe family, not absence of any useful motive.
A fair size-graded union enlarges coverage without making a fixed finite family
complete for arbitrary accumulators or invariants.

\subsection{Typed transports and restricted canonicalization}\label{sec:dependent-transports}

For $F:I\to\Sort\,u$, $v:F(i)$, and $e:i=j$, transport constructs a term of
$F(j)$, for example $\mathsf{cast}(\mathsf{congrArg}\ F\ e,v)$. Keep the family,
source and destination indices/types, equality proof, transported term, and
normalization provenance. A deferred transport must be a typed IR node whose
interpretation inserts that cast, never an ordinary Lean term temporarily
treated as though it had a different type.\prov{N1--N9}

Three notions must remain separate. Definitional equality permits conversion;
propositional equality provides transport evidence; finite observational
agreement generally does not permit dependent substitution. Proof irrelevance
identifies proofs of the same proposition, not arbitrary endpoint types or
data in different fibers. Runtime proof erasure is a fourth implementation
fact, not a license to erase the typing role of \code{Eq.rec}.

The original example needs correction: with the standard Lean natural addition,
$n+0$ reduces to $n$, whereas $0+n=n$ at symbolic $n$ uses the corresponding
theorem. Thus $\mathsf{Vec}\ A\ (0+n)$, associative index rearrangements, or
$\mathsf{Vec}\ A\ (n+m)$ versus $\mathsf{Vec}\ A\ (m+n)$ exercise propositional
transport; $\mathsf{Vec}\ A\ (n+0)$ normally does not.\cite{leannat}
Test native definitional equality before adding a cast.\prov{N1,N3--N7,N9}

Choose a terminating deterministic index normalizer $N$ for a declared
fragment, returning $e_i:i=N(i)$. Constructor arithmetic or canonical affine
expressions are possible first theories. Normalize endpoints, compose their
certificates, and reuse one transport construction for the normalized pair.
For vector append, one proved index equation can serve many corresponding
terms. Outside the fragment, keep a symbolic equality obligation; different
normal forms do not prove inequality unless the fragment procedure is complete
for that claim. Arbitrary \code{simp} success is a proof service, not a definition
of global canonical form.\prov{N1--N9}

Normalize identity casts and compatible compositions with checked lemmas.
Transport along $p:A=B$ followed by $q:B=C$ agrees with transport along their
composite, by equality elimination on $p$ and $q$. This optimization retains
endpoints. Orient rewrites or bound saturation so arithmetic equalities do not
create rewriting loops. The policy, lemma set, and certificates belong in
cache identity and receipts. Place casts consistently at constructor/provider
boundaries rather than generating every distribution through internal nodes.
Initially restrict registered families and index theories.\prov{N4,N5,N8,N9}

Heterogeneous equality can express relations before types coincide; converting
that relation into the final homogeneous inhabitant still needs reconstruction.
Typed e-graphs or local congruence classes may share checked equalities, but
dependency-aware congruence and transport-producing extraction are obligations.
A setoid changes the semantic language and requires every operation and
contract to respect its relation. Neither provides a general cheap quotient of
all Lean terms by arbitrary provable equality.

Proof evidence may have low or zero presentation/ranking weight, but transport
work remains real. Use a positive syntax grade, a certified canonical
representation, or an explicit finite transport bound to avoid infinitely
many zero-cost paths. The earlier bound of two casts is an experiment that
restricts the grammar, not a completeness theorem. Hard deduplication requires
a common expected type and checked definitional equality or an equality
certificate; a coarser proof-insensitive fingerprint may group, not delete,
unequivalent candidates. If display hides casts, either emit the exact checked
term too or re-elaborate the sugared source and certify its intended equality.
Do not publish untypeable Lean as a paste-ready result.\prov{N1--N9}

\subsection{Demand-driven unfolding and index narrowing}\label{sec:dependent-unfolding}

Canonical describes exposing function aliases such as \code{Set} and
\code{Not}, reducible definitions, and class-instance structure in its
translation; it also notes dependency expansion and abstraction costs.
The documented \code{sauto} controls distinguish heuristic \code{unfold:}
(default explicit list empty) from forced \code{unfold!:}; basic logical
definitions receive special treatment. \code{red} simplification and
\code{ered} eager scheduling are documented default-on options, with separate
context-change simplification and backtracking solvers. N6 explicitly did not
verify those implementation defaults; the more precise statements come from
the documentation evidence in N1, N3, N7, and N9, not from a comparative run.
Pin revisions and options in any experiment.\cite{canonical,sautodocs}

Use a phase-sensitive transparency ladder:
\begin{enumerate}
\item expose only the binder, function spine, constructor family, class head,
or rigid application head needed for retrieval or the next rule;
\item normalize registered index expressions using checked equations;
\item unfold a demanded head blocking matching, a scrutinee, or motive
construction, charging the particular request;
\item offer a bounded broader-unfolding or theorem-rewrite fallback.
\end{enumerate}
Retrieval keeps cheap normalized heads; speculative unification can escalate
for a concrete proposed application; observation evaluation uses supported
definitions/equations; proof simplification uses a bounded explicit rewrite
set. Ordinary provider names remain compact components even when their bodies
are unfolded for reasoning. Expanding the whole transitive provider closure
can destroy indexing selectivity and duplicate enormous expressions.\prov{N1--N9}

Record the request's purpose, transparency, unfolded declarations, environment,
local context, and assignment dependencies. A shallow mismatch means unresolved
at this stage, not proof of nonconvertibility. It cannot permanently remove a
candidate requiring a later stage allowed by the advertised grammar. Escalation
is resumable work charged to the ledger, with per-constant budgets and cycle
protection. Opaque constants remain opaque unless an allowed checked equation
or explicit interface permits reasoning about them.

If $F(?n)$ is stuck because the index determines its head, repeatedly unfolding
$F$ adds no information. Schedule another goal that determines $?n$, or
enumerate constructor forms of the inductive index under the split budget.
This is \emph{narrowing}, a search alternative sharing the native assignment
state, not a normalization theorem. Keep all admitted forms and their
dependencies.\prov{N7}

Compare fixed weak-head/reducible policies, demand-driven escalation, and
bounded up-front normalization on indexed vectors, arithmetic, and large
libraries. Count matches unlocked, failed conversions, reduction work,
expression growth, cache invalidations/memory, and verified candidates at
matched logical budgets. Compare Canonical, \code{exact?}, and translated
\code{sauto}/Agda cases to locate capability boundaries; changed elaboration
and libraries make raw cross-assistant timings informative case studies,
not an implementation theorem or a demonstrated optimal policy.\prov{N1--N9}

\subsection{A version-pinned native motive adapter}\label{sec:dependent-adapter}

The Lean~4.34.0 source inspected by the responses exposes
\code{Lean.Meta.getMajorTypeIndices},
\code{Lean.Meta.mkRecursorAppPrefix}, and \code{Lean.MVarId.induction}.
The first requires suitable variable indices and checks repeated or improperly
dependent occurrences. The prefix builder abstracts the major premise and
indices, chooses recursor universes, and checks the motive/elimination sort.
Induction reverts dependencies, reintroduces them, and returns
\code{InductionSubgoal} values containing goals, fields, and a
\code{FVarSubst}.\cite{leaninduction}\prov{N1--N9}

\begin{lstlisting}[caption={Signatures reported from the pinned source; integration still needs compilation tests.}]
Lean.Meta.mkRecursorAppPrefix
  (mvarId : MVarId) (tacticName : Name)
  (majorFVarId : FVarId) (recursorInfo : RecursorInfo)
  (indices : Array Expr) : MetaM Expr

Lean.MVarId.induction
  (mvarId : MVarId) (majorFVarId : FVarId)
  (recursorName : Name)
  (givenNames : Array AltVarNames := #[]) :
  MetaM (Array InductionSubgoal)
\end{lstlisting}
Do not duplicate reversion already performed by the native API, and do not
assume success of high-level \code{cases} means every low-level induction call
accepts the unchanged major premise. Compound/repeated indices may need the
elaborator's generalization path or remembered equalities first. N3 also
identifies \code{ElimApp.mkElimApp} and \code{ElimApp.setMotiveArg} in tactic
elaboration as implementation references, with their surrounding state and
preconditions intact. Functional induction over an existing definition is
another available principle; it does not itself invent the desired recipe.\cite{leanfuninduction,lean418}

A trial accepts the major premise, generalized telescope, recursor/principle,
and recipe. Save complete speculative state, validate the recipe, invoke the
native operation, inspect branch targets, and retain substitutions and
transports. Return a persistent successor or replay recipe; raw metavariable
identifiers from a rolled-back state must not escape. Apply the same
substitution to observations and recursive capabilities instead of guessing
new locals from printed names or positions.\prov{N1--N9}

Use structured outcomes:
\begin{lstlisting}[caption={Proposed adapter results; pseudocode.}]
MotiveAttempt :=
  success(prefix, branches, substitutions, transports, replay)
| unsupportedMajorOrRecursor(reason)
| indexPreparationNeeded(indices)
| illegalEliminationOrUniverse(reason)
| illTypedMotive(diagnostic)
| missingInstanceOrUnresolvedConstraint(constraint)
| resourceExhausted
\end{lstlisting}
Preserve native diagnostics as detail, not fragile exception-text predicates
for impossibility. Failed attempts restore assignments, universes, local
instances, messages, delayed constraints, substitutions, and branch caches.
Cancellation and internal/resource interruption propagate to the controller;
they are not ordinary motive failures. Limit heartbeats/recursion per trial,
but do not promise cheap failure merely because the function is callable from
\code{MetaM}. Memoize only with the complete goal/context, recipe, recursor,
policy, and relevant assignment identity.

The first adapter suite should include lists, naturals, options, vector map,
vector zip, dependent pairs, \code{Fin n} locals, compound/repeated indices,
major-value-dependent result families, forbidden elimination from \code{Prop},
and a near-identical ill-scoped recipe. Repeated failed trials must leave the
next trial observationally unchanged. Measure instantiation and failure cost
separately from branch search. Compile against the actual pinned toolchain
before declaring this a supported API, and keep specialized fast constructors
behind the same validation boundary.\prov{N1--N9}

\subsection{New questions after the four original answers}\label{sec:dependent-open}

The finite-template, transport, unfolding, and API-access questions above have
concrete bounded answers. Their unresolved extensions are:
\begin{enumerate}[label=R4.N\arabic*.]
\item Which explicit motive recipe grammar, including nonmaximal index
abstractions and dependency-closed changing parameters, gives the best measured
coverage before new carrier or invariant invention is required?
\item For which registered dependent families and index theories can transport
normalization be proved terminating and reconstructible while keeping its
enumeration grade and proof cost bounded?
\item Which phase-sensitive transparency policy improves verified results at
matched logical budgets without turning shallow conversion failures into
permanent exclusions?
\item Can the pinned induction adapter preserve substitutions and every
speculative state component across failure and interruption, and what are its
measured success/failure costs on compound indices and value-dependent motives?
\end{enumerate}
